\documentclass[11pt]{amsart}
\usepackage{calc,amssymb,amsmath,ulem,stmaryrd,fullpage}
\usepackage{enumerate}
\usepackage{alltt}
\RequirePackage[dvipsnames,usenames]{color}
\let\mffont=\sf
\normalem
%\input{kmacros3.sty}
\input{xy}
\xyoption{all}
\usepackage{hyperref}
\usepackage{graphicx}
\usepackage[all,cmtip]{xy}
\usepackage{verbatim}
\usepackage[left=0.95in,top=0.95in,right=0.95in,bottom=0.95in]{geometry}
\newcommand{\tauCohomology}{T}
\newcommand{\FCohomology}{S}


\theoremstyle{theorem}
%\newtheorem*{mainthm*}{Main Theorem}

\newcommand{\note}[1]{\marginpar{\sffamily\tiny #1}}

\def\todo#1{\textcolor{Mahogany}%
{\sffamily\footnotesize\newline{\color{Mahogany}\fbox{\parbox{\textwidth-15pt}{#1}}}\newline}}

\renewcommand{\O}{\mathcal O}
\newcommand{\ie}{{\itshape i.e.} }
\newcommand{\roundup}[1]{\lceil #1 \rceil}
\newcommand{\rounddown}[1]{\lfloor #1 \rfloor}
\renewcommand{\to}[1][]{\xrightarrow{\ #1\ }}
\newcommand{\onto}[1][]{\protect{\xrightarrow{\ #1\ }\hspace{-0.8em}\rightarrow}}
\newcommand{\into}[1][]{\lhook \joinrel \xrightarrow{\ #1\ }}
%\newcommand{DBDual}[1]{{\underline{\omega}_{#1}}}
\newcommand{\DBDual}[1]{{{\uuline \omega}{}^{\mydot}_{#1}}}
\newcommand{\Tt}{{\mathfrak{T}}}
%\newcommand{\bn}{{\mathfrak{n}} }
\newcommand{\sol}[1]{ \vskip 6pt{\color{blue} {\bf Solution:} #1} \vskip 6pt}

\newcommand{\bR}{{\mathbb{R}}}
\newcommand{\va}{{\vec{a}}}
\newcommand{\vb}{{\vec{b}}}
\newcommand{\vzero}{{\vec{0}}}
\newcommand{\vi}{{\vec{i}}}
\newcommand{\vj}{{\vec{j}}}
\newcommand{\vk}{{\vec{k}}}
\newcommand{\vh}{{\vec{h}}}
\newcommand{\vr}{{\vec{r}}}
\newcommand{\vu}{{\vec{u}}}
\newcommand{\vv}{{\vec{v}}}
\newcommand{\vw}{{\vec{w}}}
\newcommand{\vx}{{\vec{x}}}
\newcommand{\vy}{{\vec{y}}}
\newcommand{\vz}{{\vec{z}}}
\newcommand{\vT}{{\vec{T}}}
\newcommand{\vN}{{\vec{N}}}
\newcommand{\vF}{{\vec{F}}}
\newcommand{\vG}{{\vec{G}}}
\newcommand{\bF}{\mathbb{F}}
\newcommand{\bP}{\mathbb{P}}
\newcommand{\bQ}{\mathbb{Q}}
\newcommand{\bZ}{\mathbb{Z}}
\newcommand{\bA}{\mathbb{A}}
\newcommand{\codim}{{\mathrm{codim}}}
\newcommand{\Map}{{\mathrm{Map}}}


\begin{document}
\title{Worksheet \#5 -- Math 6130\\Fall 2018}
\author{Due Monday, October 15th}

\maketitle
You may work in groups of up to 3.  Only one worksheet needs to be turned in per group.
\vskip 9pt
\noindent
{\bf 1.}  For any two points $p = (a_0 : a_1 : \dots : a_n), q = (b_0 : b_1 : \dots :b_n) \in \bP^n$, show that there is a unique projective line $L$ through those two points and compute the ideal $I(L)$ (based on the above coordinates).  
\vskip 200pt
\noindent
{\bf 2.}  Consider $H \subseteq \bP^n$ defined by $Z(x_n)$.  For a point $(a_0 : a_1 : \dots : a_n = 1) = p \in \bP^n \setminus H$, we can create a map 
\[
U = \bP^n \setminus \{ p \} \xrightarrow{\phi_p} H \cong \bP^{n-1}
\]
which takes $q \in U$, considers the line $L$ through $p$ and $q$, and takes $q$ to the point $H \cap L$.  Write down formulas for the coordinates of this map and verify it is regular on $U$.  
\newpage
\noindent
{\bf 4.}  Suppose that $E$ is a $d$-dimensional linear subspace of $\bP^n$.  Let $L_1, \dots, L_{n-d}$ be linear forms in $k[x_0, \dots, x_n]$ which define $E$.  We consider the rational map
\[
\phi : P^n \xrightarrow{(L_1: \dots : L_n)} \bP^{n-d-1}.
\]
Choose a linear subspace $F$ of $\bP^n$ of dimension $n-d-1$ which is disjoint from $E$.  (Simply extend $L_1, \dots, L_{n-d}$ to a $k$-basis of the 1-degree elements of $S(\bP^n)$, say $L_1, \dots, L_{n-d}, M_1, \dots M_{d+1}$ and let those forms define $F$).
For each $p \in \bP^n \setminus E$, construct $G_p$ the unique $(d+1)$-dimensional linear subspace of $\bP^n$ containing $p$ and $E$.  Show that $G_p$ intersects $F$ in a unique point.
\vskip 230pt
\noindent
{\bf 3.}  With notation as in {\bf 2.}, consider $W = Z(x_n, x_{n-1})$ and choose a point $q \in H \setminus W$.  Again construct a regular map $V = H \setminus W \xrightarrow{\phi_q} W$ as in {\bf 2.}  Now, compute the natural domain of the rational map
\[
\phi_q \circ \phi_p
\]
and describe the map geometrically as a projection from a linear subspace in the sense of the previous problem.
\end{document}

